\documentclass[11pt,a4paper]{article}
\usepackage{fontspec,amsmath,amssymb,amsthm,microtype}
\setmainfont{lmroman10-regular.otf}[BoldFont=lmroman10-bold.otf,ItalicFont=lmroman10-italic.otf,BoldItalicFont=lmroman10-bolditalic.otf]
\setmonofont{lmmono10-regular.otf}[Scale=MatchLowercase]
\usepackage[margin=25mm]{geometry}
\usepackage[unicode,hidelinks]{hyperref}
\setlength{\emergencystretch}{3em}
\setlength{\parskip}{0.4em}
\setlength{\parindent}{0pt}
\allowdisplaybreaks[2]
\raggedbottom
\title{\texorpdfstring{Spectral cancellation and pathwise exponential decay}{Spectral cancellation and pathwise exponential decay}}
\author{LangevinLean mathematical exposition}
\date{}
\begin{document}\maketitle
\noindent Source module: \texttt{Kinetic\allowbreak{}Langevin/\allowbreak{}Contraction.\allowbreak{}lean}. The accompanying HTML pairs this mathematical exposition with the exact Lean statements and records the scope of the formalization.


This module proves a deterministic theorem for a kinetic difference equation \(v'=-\gamma v-H_tx\), \(x'=v\). The continuous linear operator \(H_t\) is allowed to vary with time but must be self-adjoint and satisfy quadratic-form bounds \(mI\preceq H_t\preceq MI\). The theorem derives energy decay from this differential equation. Constructing stochastic solutions whose difference has these properties and taking an infimum over probability couplings are separate layers.


\section{\texorpdfstring{Write the derivative and prove exact spectral cancellation}{Write the derivative and prove exact spectral cancellation}}


The transformed derivatives are \(y'=-av-Hx\), \(z'=(\gamma-a)v+Hx\). Thus \(\mathcal E'=2\langle y,y'\rangle+2\langle z,z'\rangle\), which \texttt{energy\allowbreak{}Derivative} records. Multiplying by \(d=\gamma-2a\) gives

\[d\mathcal E'=2\bigl(a^2\|y\|^2-\langle y,Hy\rangle+\langle z,Hz\rangle-(\gamma-a)^2\|z\|^2\bigr).\]

The proof expands linear-map applications, scalar multiplication and inner products. Self-adjointness rewrites \(\langle v,Hx\rangle\) as \(\langle x,Hv\rangle\), so cross terms cancel. \texttt{ring} verifies the remaining scalar polynomial identity. Clearing the denominator keeps the identity valid even at singular parameter values.



{\raggedright\noindent Lean declarations: \texttt{energy\allowbreak{}Derivative}, \texttt{spectral\_\allowbreak{}identity}.\par}

\section{\texorpdfstring{Use lower and upper curvature on different coordinates}{Use lower and upper curvature on different coordinates}}


Apply the lower bound \(\langle y,Hy\rangle\ge m\|y\|^2\) and upper bound \(\langle z,Hz\rangle\le M\|z\|^2\). Let \(k=\max\{a^2-m,M-(\gamma-a)^2\}\). Both coefficients are at most \(k\); multiply these scalar bounds by the nonnegative coordinate squares. The spectral identity then yields \(d\mathcal E'\le2k\mathcal E\). The explicit premise \(d>0\) permits division without reversing the inequality. Lean names each quadratic-form and maximum comparison, combines them with \texttt{nlinarith}, then uses \texttt{le\_\allowbreak{}div\_\allowbreak{}iff\ensuremath{{}_{0}}}. This proves the dissipativity estimate from spectral bounds rather than assuming contraction.



{\raggedright\noindent Lean declarations: \texttt{spectral\_\allowbreak{}dissipativity}.\par}

\section{\texorpdfstring{Specialize to the simple friction regime}{Specialize to the simple friction regime}}


At \(a=0\), the maximum becomes \(\max\{-m,M-\gamma^2\}\). The assumption \(m+M\le\gamma^2\) gives \(M-\gamma^2\le-m\), so the maximum equals \(-m\). With \(\gamma>0\), the differential bound becomes \(\mathcal E'\le-2m\mathcal E/\gamma\). Lean obtains the general bound, proves the scalar comparison by \texttt{linarith}, and simplifies the maximum. This is the energy exponent; taking a square root later gives the norm exponent \(-m/\gamma\).



{\raggedright\noindent Lean declarations: \texttt{spectral\_\allowbreak{}dissipativity\_\allowbreak{}zero}.\par}

\section{\texorpdfstring{Differentiate the actual transformed energy along a path}{Differentiate the actual transformed energy along a path}}


Assume \texttt{Has\allowbreak{}Deriv\allowbreak{}At} witnesses the two kinetic equations at time \(t\). The linear combinations \(v+(\gamma-a)x\) and \(-v-ax\) inherit derivatives by addition, negation and constant scalar multiplication. The vector normalization tactic \texttt{module} reduces them to \(-av-Hx\) and \((\gamma-a)v+Hx\). The Hilbert-space squared-norm chain rule gives derivative \(2\langle y,y'\rangle\), and adding the two rules produces \texttt{energy\allowbreak{}Derivative}. No derivative of \(H_t\) is needed: only the value at the current time enters the first-order kinetic equation.



{\raggedright\noindent Lean declarations: \texttt{has\allowbreak{}Deriv\allowbreak{}At\_\allowbreak{}transform\allowbreak{}Energy}.\par}

\section{\texorpdfstring{Apply Gronwall with an explicitly proved derivative bound}{Apply Gronwall with an explicitly proved derivative bound}}


Define scalar functions \(F(t)=\mathcal E(v_t,x_t)\) and \(F'(t)=\texttt{energyDerivative}(H_t,v_t,x_t)\), with \(K=2k/(\gamma-2a)\). Continuous paths make \(F\) continuous on the closed horizon, and the preceding derivative theorem gives its derivative on the half-open interval. Spectral dissipativity proves \(F'\le KF+0\). Lean feeds these facts to Mathlib's right-slope Gronwall theorem and simplifies its zero-forcing formula, obtaining \(F(t)\le F(0)e^{Kt}\). The zero additive term is significant: this is exact synchronous-flow contraction rather than a numerical defect estimate.



{\raggedright\noindent Lean declarations: \texttt{transformed\_\allowbreak{}energy\_\allowbreak{}decay}.\par}

\section{\texorpdfstring{Project the energy estimate to the position difference}{Project the energy estimate to the position difference}}


Suppose \(v_0=0\). Substitute the exact initial energy \(((\gamma-a)^2+a^2)\|x_0\|^2\) into the Gronwall estimate. The position projection theorem gives \((\gamma-2a)^2\|x_t\|^2\le2F(t)\). Combining the two inequalities yields the stated squared version of Proposition 1. The final \texttt{nlinarith} only assembles already proved bounds. Its square root, under \(\gamma-2a>0\), is the paper's prefactor \(\sqrt{2((\gamma-a)^2+a^2)}/(\gamma-2a)\). This theorem concerns paths satisfying the explicit ODE; it is not yet a probability-law mixing result.



{\raggedright\noindent Lean declarations: \texttt{synchronous\_\allowbreak{}position\_\allowbreak{}sq\_\allowbreak{}decay}.\par}

\section{\texorpdfstring{Include the initial time through right derivatives}{Include the initial time through right derivatives}}


Finite-horizon integral solutions naturally provide a derivative within \([t,\infty)\), including at \(t=0\). The coordinate differentiation proof works unchanged with \texttt{Has\allowbreak{}Deriv\allowbreak{}Within\allowbreak{}At}: add the two right derivatives, normalize their vector coefficients, apply the squared-norm derivative rule, and sum. The resulting derivative is still \texttt{energy\allowbreak{}Derivative}. The right-slope Grönwall theorem therefore gives the same exponential energy bound under the spectral hypotheses, without imposing a differentiable extension of the stochastic paths to negative time.



{\raggedright\noindent Lean declarations: \texttt{has\allowbreak{}Deriv\allowbreak{}Within\allowbreak{}At\_\allowbreak{}transform\allowbreak{}Energy}, \texttt{transformed\_\allowbreak{}energy\_\allowbreak{}decay\_\allowbreak{}right}.\par}

\section{\texorpdfstring{Recover the full arbitrary-\allowbreak{}shift prefactor}{Recover the full arbitrary-shift prefactor}}


Write \(d=\gamma-2a>0\) and \(k=\max\{a^2-m,M-(\gamma-a)^2\}/d\). The energy estimate and equal initial velocities give \(d^2\|x_t\|^2\le2((\gamma-a)^2+a^2)e^{2kt}\|x_0\|^2\). Lean proves \((e^{kt})^2=e^{2kt}\) using \texttt{Real.\allowbreak{}exp\_\allowbreak{}add}, proves the square-root square identity from nonnegativity of its argument, and compares the squares of two nonnegative quantities. This first yields \(d\|x_t\|\le\sqrt{2((\gamma-a)^2+a^2)}e^{kt}\|x_0\|\). Dividing by the positive \(d\) gives precisely Proposition 1's prefactor and exponent.



{\raggedright\noindent Lean declarations: \texttt{synchronous\_\allowbreak{}position\_\allowbreak{}decay\_\allowbreak{}right}.\par}

\end{document}
